% ============================================================
\section{Phần I: Văn phạm, parser và sinh câu}
\label{sec:part1}
% ============================================================

\subsection{Tokenizer}

Tokenizer là bước tiền xử lý dùng chung cho parser, bộ sinh câu và entity lexicon,
nhờ đó một câu sinh ra từ văn phạm luôn được tách token giống hệt khi parse lại. Các
bước xử lý:
\begin{enumerate}
  \item Chuẩn hóa Unicode về dạng NFC và chuyển chữ thường, để ``Tuần'' gõ bằng các bộ
        gõ khác nhau (dựng sẵn hay tổ hợp dấu) cho cùng một token.
  \item Chuẩn hóa cách viết kỹ thuật: \texttt{L.O.2.5} $\rightarrow$ \texttt{lo2.5};
        bỏ dấu ngoặc kép; gộp khoảng trắng thừa.
  \item Tách dấu câu \texttt{? ! . , : ;} thành token riêng, nhưng giữ nguyên các từ
        ghép kỹ thuật như \texttt{context-free}, \texttt{llm/rag}, \texttt{lo2.5}
        và chữ viết tắt có dấu chấm như \texttt{TP.HCM}.
\end{enumerate}
Ví dụ: \texttt{``Môn NLP có bao nhiêu tín chỉ?''} $\rightarrow$
\texttt{môn / nlp / có / bao / nhiêu / tín / chỉ / ?}.

\subsection{Thiết kế văn phạm}

\paragraph{Định dạng.} Văn phạm nằm trong \code{data/grammar.cfg}: ký hiệu bắt đầu khai
báo bằng \texttt{\%start S}, mỗi luật có dạng \texttt{A -> ...}, các lựa chọn phân tách
bằng \texttt{|}, terminal đặt trong ngoặc kép và viết thường, một lựa chọn rỗng là
luật $\varepsilon$. Văn phạm hiện có 100 ký hiệu không kết thúc, 463 luật sinh (tính
theo từng lựa chọn) và 330 terminal.

\paragraph{Văn phạm ngữ nghĩa.} Nhóm thiết kế văn phạm theo hướng \emph{semantic
grammar}: tên ký hiệu không kết thúc mang nghĩa của miền thay vì nhãn từ loại. Luật gốc
tách phần lịch sự khỏi nội dung câu hỏi:
\begin{grammar}
S        -> PREFIX QUESTION SUFFIX
PREFIX   -> | PLEASE_PREFIX                 # "cho mình hỏi", "xin hỏi", ...
SUFFIX   -> | END_MARK | POLITE | POLITE END_MARK
QUESTION -> Q_COURSE_CREDITS | Q_SCHEDULE_WEEK | ... | Q_CONTEXT_LO_CHAPTER   # 26 nhánh
\end{grammar}
Mỗi câu hợp lệ đi qua đúng một trong \textbf{26 nhánh} \texttt{Q\_*} của
\texttt{QUESTION}; nhánh này quyết định intent (\cref{tab:branches}). Thực thể được bắt
bằng các thành phần có tên như \texttt{WEEK\_ENTITY}, \texttt{CHAPTER\_ENTITY},
\texttt{TOPIC\_ENTITY}, nên tầng ngữ nghĩa đọc thẳng các node này trên cây mà không phải
tìm lại từ khóa trong câu thô. Ví dụ nhánh hỏi lịch theo tuần:
\begin{grammar}
Q_SCHEDULE_WEEK       -> WEEK_ENTITY WEEK_CONTENT_QUESTION
                       | "nội" "dung" WEEK_ENTITY WEEK_CONTENT_QUESTION
                       | "nội" "dung" "của" WEEK_ENTITY
WEEK_ENTITY           -> "tuần" WEEK_NUMBER
WEEK_NUMBER           -> "1" | "01" | "2" | "02" | ... | "15" | "16" | "20" | "99"
WEEK_CONTENT_QUESTION -> "học" "gì" | "học" "những" "gì" | "có" "nội" "dung" "gì"
                       | "học" "chương" "nào" | "giảng" "chủ" "đề" "gì" | ...
\end{grammar}
\texttt{WEEK\_NUMBER} cố ý nhận cả số tuần không có trong KB (16, 20, 99): câu ``Tuần 99
học gì?'' là câu hợp lệ về cú pháp, nên hệ thống trả lời ``không tìm thấy dữ kiện''
thay vì coi là câu ngoài văn phạm.

\begin{table}[htbp]
  \centering
  \small
  \begin{tabularx}{\linewidth}{|P{3.3cm}|L|}
    \hline
    \textbf{Intent} & \textbf{Nhánh \texttt{Q\_*} trong văn phạm} \\ \hline
    \code{GET_COURSE_INFO} & \code{Q_COURSE_CREDITS}, \code{Q_COURSE_PARTS},
      \code{Q_COURSE_DESCRIPTION}, \code{Q_COURSE_EXAM_DURATION},
      \code{Q_COURSE_EXAM_TIME}, \code{Q_COURSE_ASSESSMENT} \\ \hline
    \code{GET_INSTRUCTOR} & \code{Q_INSTRUCTOR} \\ \hline
    \code{GET_SCHEDULE} & \code{Q_SCHEDULE_WEEK}, \code{Q_SCHEDULE_CHAPTER},
      \code{Q_SCHEDULE_TOPICS} \\ \hline
    \code{GET_TOPIC} & \code{Q_TOPIC_BY_NAME}, \code{Q_TOPIC_BY_CHAPTER},
      \code{Q_TOPIC_RELATION} \\ \hline
    \code{GET_LO} & \code{Q_LO} \\ \hline
    \code{GET_ASSIGNMENT} & \code{Q_ASSIGNMENT_OVERVIEW}, \code{Q_ASSIGNMENT_PART},
      \code{Q_ASSIGNMENT_FEATURE}, \code{Q_ASSIGNMENT_OUTPUT} \\ \hline
    \code{GET_DEADLINE} & \code{Q_DEADLINE} \\ \hline
    \code{GET_RESOURCE} & \code{Q_RESOURCE} \\ \hline
    \code{GET_RULE} & \code{Q_RULE} \\ \hline
    \code{CONTEXT_DEPENDENT} & \code{Q_CONTEXT_TOPIC_RELATION},
      \code{Q_CONTEXT_NEXT_CHAPTER}, \code{Q_CONTEXT_ASSIGNMENT_OUTPUT},
      \code{Q_CONTEXT_RESOURCE}, \code{Q_CONTEXT_LO_CHAPTER} \\ \hline
    \code{UNKNOWN} & Không có cây (câu ngoài văn phạm) \\ \hline
  \end{tabularx}
  \caption{26 nhánh câu hỏi và intent tương ứng}
  \label{tab:branches}
\end{table}

\paragraph{Quá trình xây dựng.} Ở giai đoạn 1, mỗi thành viên viết các nhánh thuộc phần
mình phụ trách (\cref{tab:members}); bản gộp giữ nguyên danh sách 26 nhánh đã chốt, chỉ
nhận các lựa chọn mới vào nhánh có sẵn, dùng lại các biến chung (\texttt{YES\_NO},
\texttt{TOPIC\_ENTITY}, \texttt{ASSIGNMENT\_ENTITY}, \texttt{CONTEXT\_REFERENCE},
\dots) và loại các lựa chọn gây mơ hồ hoặc làm rớt câu đang đúng. Sau giai đoạn 1, nhóm
bổ sung nhóm thành phần tính năng của bài tập lớn: các cách gọi một tính năng được gom
vào node mang tên phần tương ứng (\cref{tab:feat}), nhờ vậy tầng ngữ nghĩa xác định phần
BTL bằng nhãn trên cây thay cho danh sách từ khóa riêng.
\begin{grammar}
ASSIGNMENT_FEATURE -> FEAT_PARSING | FEAT_SEMANTIC | FEAT_CONTEXT | FEAT_EVAL
FEAT_PARSING       -> "parser" | "grammar" | "parse" "tree" | "văn" "phạm"
                    | "phân" "tích" "cú" "pháp"
\end{grammar}

\begin{table}[htbp]
  \centering
  \small
  \begin{tabular}{|l|c|l|}
    \hline
    \textbf{Node} & \textbf{Phần BTL} & \textbf{Ví dụ cách gọi} \\ \hline
    \code{FEAT_PARSING} & I & parser, grammar, parse tree, phân tích cú pháp \\ \hline
    \code{FEAT_SEMANTIC} & II & semantic, hỏi đáp, intent, entity, KB \\ \hline
    \code{FEAT_CONTEXT} & III & ngữ cảnh, hội thoại, tham chiếu \\ \hline
    \code{FEAT_EVAL} & IV & đánh giá, test set, metric, kiểm thử \\ \hline
  \end{tabular}
  \caption{Nhóm thành phần tính năng của bài tập lớn}
  \label{tab:feat}
\end{table}

\subsection{Parser Earley}

\paragraph{Lựa chọn thuật toán.} Nhóm dùng thuật toán Earley \cite{earley1970,jurafsky}
vì nó xử lý được mọi CFG, kể cả luật $\varepsilon$ (\texttt{PREFIX}, \texttt{SUFFIX}) và
luật đệ quy trái, không cần đưa văn phạm về dạng chuẩn Chomsky như CKY, nên cây trả về
giữ đúng các node ngữ nghĩa mà tầng sau cần đọc. Độ phức tạp trong trường hợp xấu nhất
là $O(n^3)$ theo số token; câu hỏi trong miền chỉ vài chục token nên chi phí không đáng kể.

\paragraph{Cách hoạt động.} Với $n$ token, parser duy trì các tập trạng thái
(chart) $S_0, \dots, S_n$. Mỗi trạng thái gồm một luật, vị trí dấu chấm trong vế phải,
vị trí bắt đầu và các cây con đã dựng. Ba thao tác lặp lại cho đến khi tập trạng thái
không đổi:
\begin{itemize}
  \item \textbf{Predictor}: khi ký hiệu sau dấu chấm là ký hiệu không kết thúc $B$,
        thêm mọi luật $B \rightarrow \bullet\,\gamma$ bắt đầu tại vị trí hiện tại.
  \item \textbf{Scanner}: khi ký hiệu sau dấu chấm là terminal trùng với token kế tiếp,
        đưa trạng thái đã dịch dấu chấm sang tập $S_{i+1}$.
  \item \textbf{Completer}: khi một trạng thái $B \rightarrow \gamma\,\bullet$ hoàn tất,
        dịch dấu chấm cho mọi trạng thái đang chờ $B$ tại vị trí bắt đầu của nó và gắn
        node $B$ vừa dựng làm cây con.
\end{itemize}
Câu được chấp nhận khi và chỉ khi $S_n$ có trạng thái \texttt{S} hoàn tất bắt đầu từ vị
trí 0, tức là parser đã tiêu thụ hết token. Vì vậy ``Tuần 3 học gì? xyz'' bị từ chối do
còn token thừa, và câu bị từ chối được xuất dưới dạng \texttt{()}.

\paragraph{Luật $\varepsilon$.} Một node $\varepsilon$ hoàn tất ngay tại vị trí nó bắt
đầu, nên completer có thể đã chạy trước khi một trạng thái chờ ký hiệu đó xuất hiện. Để
không bỏ sót dẫn xuất, parser giữ chỉ mục các node rỗng đã hoàn tất theo vị trí; khi
predictor gặp ký hiệu rỗng được, nó lấy node từ chỉ mục này và dịch dấu chấm ngay.

\paragraph{Phát hiện mơ hồ.} \code{parse()} trả về cây đầu tiên theo thứ tự luật;
\code{parse_trees(limit=2)} tìm tối đa hai cây khác nhau, đủ để kết luận một câu có mơ
hồ hay không. Số dẫn xuất lưu cho mỗi mục trong chart cũng bị giới hạn để tránh lặp vô
hạn khi có chu trình qua luật rỗng. Đây là cơ chế phát hiện mơ hồ có giới hạn, không
phải thuật toán liệt kê mọi cây (parse forest).

\paragraph{Ví dụ.} \Cref{fig:tree} là cây cú pháp của câu ``Tuần 3 học gì?'', dạng
ngoặc tương ứng trong \code{parse-results.txt}:
\begin{trace}
(S (PREFIX ε) (QUESTION (Q_SCHEDULE_WEEK (WEEK_ENTITY tuần (WEEK_NUMBER 3))
   (WEEK_CONTENT_QUESTION học gì))) (SUFFIX (END_MARK ?)))
\end{trace}

\begin{figure}[H]
  \centering
  \begin{adjustbox}{max width=\linewidth}
  \begin{forest}
    for tree={font=\footnotesize\ttfamily, s sep=4mm, l sep=4mm, inner sep=1.5pt,
      edge={black!60}},
    where n children=0{font=\small\itshape, text=blue!50!black}{}
    [S
      [PREFIX [$\varepsilon$]]
      [QUESTION
        [Q\_SCHEDULE\_WEEK
          [WEEK\_ENTITY [tuần] [WEEK\_NUMBER [3]]]
          [WEEK\_CONTENT\_QUESTION [học] [gì]]]]
      [SUFFIX [END\_MARK [?]]]]
  \end{forest}
  \end{adjustbox}
  \caption{Cây cú pháp của câu ``Tuần 3 học gì?''}
  \label{fig:tree}
\end{figure}

Câu ngoài miền như ``Hôm nay trời có mưa không?'' không có dẫn xuất nào, được xuất là
\texttt{()} và nhận intent \code{UNKNOWN}; hệ thống không truy vấn KB cho câu này.

\subsection{Sinh câu từ văn phạm}

Bộ sinh câu (\code{generator.py}) dẫn xuất từ \texttt{S}, luôn mở rộng ký hiệu không kết
thúc \emph{trái nhất} của dạng câu hiện tại:
\begin{itemize}
  \item Mỗi nhánh \texttt{Q\_*} có một hàng đợi dẫn xuất riêng (duyệt theo chiều rộng);
        các hàng đợi được lấy xoay vòng để mọi loại câu hỏi đều có mặt trong kết quả
        thay vì dồn vào vài nhánh đầu.
  \item Trước khi sinh, công cụ tính số token tối thiểu mà mỗi ký hiệu có thể sinh ra;
        dạng câu nào chắc chắn vượt giới hạn 32 token bị loại sớm. Cùng với giới hạn độ
        dài dạng câu, điều này bảo đảm quá trình luôn dừng kể cả với luật đệ quy và luật
        rỗng; câu trùng bị loại.
  \item Câu được hoàn thiện theo hậu tố mà văn phạm cho phép, viết hoa chữ đầu và giữ
        nguyên chữ viết tắt kỹ thuật.
\end{itemize}
Trước khi ghi \code{samples.txt}, CLI \textbf{parse lại mọi câu sinh ra}; câu nào không
parse được sẽ làm lệnh thất bại. File có tối đa 10.000 câu theo yêu cầu đề; lần chạy mặc
định sinh 1.000 câu, ví dụ:
\begin{trace}
Nội dung này nằm ở chương nào?
Chương tiếp theo thì sao?
Tôi cần hoàn thành những bài tập nào?
\end{trace}

\subsection{Kiểm tra văn phạm}
\label{sec:grammarcheck}

Công cụ \code{tools/check_grammar.py} kiểm tra tĩnh và động văn phạm sau mỗi lần sửa:
\begin{itemize}
  \item \textbf{Tĩnh}: định nghĩa trùng, luật trùng, ký hiệu không kết thúc chưa định
        nghĩa hoặc không tới được từ \texttt{S}, terminal viết hoa, terminal mà tokenizer
        sẽ tách khác đi, và đủ 26 nhánh \texttt{Q\_*}.
  \item \textbf{Động}: chạy câu mẫu, câu thử thách, hội thoại; sinh 1.000 câu rồi parse
        lại, đồng thời chạy tiếp tầng ngữ nghĩa và truy vấn trên các câu sinh.
\end{itemize}
Ở lần kiểm tra gần nhất (lưu tại \code{docs/validation/grammar-check.json}): 1.000 câu
sinh phủ đủ 26/26 nhánh (38--39 câu mỗi nhánh), tất cả parse lại được, không câu nào cho
\code{UNKNOWN} hay entity rỗng ngoài các nhánh ngữ cảnh; 1.047 câu được kiểm tra mơ hồ
và không câu nào có hai cây. Kết quả này chỉ áp dụng cho tập câu đã chạy, không phải
chứng minh văn phạm không mơ hồ.
